% ID: SEL-D0071-OB-E03
% ESTADO: cotejada. Segunda pasada visual y borrador compilado revisados.
% Fuente: Selectividad/sel_2009_matematicas_aplicadas/mat apli CC SS05.pdf
% SHA-256: ca51ab435c7584a8e047e19d0826f3273f3e88aed6110ccd00e46a40e85518b1
% Página: 2; opción B; ejercicio original 3.

\noindent\underline{Parte I}

$A$ y $B$ son dos sucesos independientes de un mismo experimento aleatorio, tales que
\[
P(A)=0.4,\quad P(B)=0.6.
\]
\begin{enumerate}[label=\alph*)]
\item \textbf{(1 punto)} Calcule $P(A\cap B)$ y $P(A\cup B)$.
\item \textbf{(1 punto)} Calcule $P(A/B)$ y $P(B/A^C)$.
\end{enumerate}
\par\noindent\begin{minipage}{\linewidth}
\noindent\underline{Parte II}

\textbf{(2 puntos)} Tomando, al azar, una muestra de 80 empleados de una empresa, se encontró que 20 usaban gafas. Halle, con un nivel de confianza del $90\%$, un intervalo de confianza para estimar la proporción de empleados de esa empresa que usan gafas.
\end{minipage}\par
